\section{Results}
\label{sec:results}

\subsection{Impulse Responses to a Contractionary Monetary Shock}
\label{sec:results_irf}

Figure~\ref{fig:irf} reports the dynamic response of output, inflation,
and interest rates to a 25 basis-point contractionary innovation in the
monetary policy rule, $\varepsilon_{t}^{m}=0.0025$ in quarterly units.
The shock follows an AR(1) process with persistence $\rho_{m}=0.50$, so
its half-life is one quarter; the propagation in the macro aggregates,
however, persists somewhat longer because of the model's forward-looking
dynamics. Three quantitative features merit discussion.

\begin{figure}[t]
  \centering
  \includegraphics[width=0.98\textwidth]{nk_irf_eps_m.pdf}
  \caption{Impulse responses to a positive 25-basis-point monetary
    policy innovation $\varepsilon_{t}^{m}$ with AR(1) persistence
    $\rho_{m}=0.50$. Output is expressed as the percentage deviation
    from steady state; inflation and interest rates are expressed as
    percentage-point deviations on an annualised basis.}
  \label{fig:irf}
\end{figure}

\paragraph{Output contracts on impact and recovers smoothly.}
Output falls by $0.259\%$ on impact and gradually returns to its steady
state, becoming indistinguishable from zero by roughly the tenth
quarter. The decline reflects the textbook New Keynesian intertemporal
substitution mechanism: the contractionary innovation raises the real
interest rate ex ante, encouraging households to defer consumption.
Because $C_{t}=Y_{t}$ by goods-market clearing, lower consumption
translates one-for-one into lower output. The magnitude of the impact
response is governed by the elasticity $1/\sigma$ in the IS
curve~\eqref{eq:nk-is}; with the textbook unit value of $\sigma$, an
$x$-basis-point increase in the real rate maps approximately into an
$x$-basis-point reduction in output, modulated by the Taylor-rule
feedback through inflation.

\paragraph{Inflation falls more than output in annualised terms.}
Inflation falls by $0.352$ percentage points on impact when expressed
on an annualised basis, and recovers along a similar timescale to
output. The Phillips-curve slope $\kappa\approx 0.172$ implies that a
$0.259$-percent decline in output translates into a quarterly inflation
decline of approximately $\kappa\,y_{t}\approx 0.044$ percent. The
annualised value of $0.178$ percentage points (i.e., $4\,\kappa\,y_{t}$)
is mechanically smaller than the full impact response, however, because
the forward-looking term in the Phillips curve also contracts:
households and firms expect future inflation to be below the steady
state, and this expected disinflation is itself a contemporaneous
Phillips-curve force pulling current inflation down further. The total annualised impact response
of $0.352$ percentage points reflects the sum of these two effects.

\paragraph{The real interest rate rises by more than the nominal rate.}
The annualised nominal rate rises by $0.342$ percentage points, while
the annualised \textit{real} rate (defined as $r_{t}-\mathbb{E}_{t}
\pi_{t+1}$) rises by $0.518$ percentage points. The real rate exceeds
the nominal rate response because expected inflation is negative
following the shock; subtracting a negative number adds to the
ex-ante real rate. This is the proximate channel through which monetary
policy contracts aggregate demand, and the asymmetric response of
real versus nominal rates is a defining feature of the New Keynesian
transmission mechanism (cf.\ \citealp{Gali2008}, Section~3.4).

\paragraph{Comovement and ranking.}
All three series respond in the directions predicted by the canonical
model: output and inflation fall together, the policy rate rises, and
the real rate rises by even more. The ordering of volatilities--with
the annualised real rate response exceeding the annualised nominal
rate response, which in turn exceeds the annualised inflation
response--is consistent with the empirical patterns documented by
\citet{ChristianoEichenbaumEvans2005} and replicates the qualitative
ranking shown in \citet[][Figure~3.1]{Gali2008}.
